\documentclass[paper=screen,mode=present,style=zysimple]{powerdot}

\usepackage{fancyvrb,verbatim,listings}
\usepackage{proof}
\usepackage[all]{xypic,xy}

\usepackage{latexsym,stmaryrd,amsmath,amsfonts}

%% \newtheorem{definition}{Definition}
%% \newtheorem{theorem}{Theorem}
%% \newtheorem{lemma}{Lemma}
\newtheorem{prop}{Proposition}
%% \newtheorem{corollary}{Corollary}
%% \newenvironment{proof}{{\sc Proof . }}{{\hfill $\square$}}
%% \newcommand{\myqed}{\hfill $\square$}
%% \newcommand{\ped}{\vspace{-1.6em}}
%% \newcommand{\note}[1]{\noindent\fbox{\parbox{\textwidth}{\em #1}}}

\newcommand{\abind}[2]{\mathtt{bind} \ (\, {#1} \,)  \ \mathtt{in} \ {#2}}
\newcommand{\acase}{\mathtt{case}}
\newcommand{\adv}{\mathcal{A}}
\newcommand{\aspecto}[1]{\xrightarrow{\; {#1} \;}}
\newcommand{\asplinear}{\mathfrak{l}}
\newcommand{\aspmodal}{\mathfrak{m}}
\newcommand{\aspnormal}{\mathfrak{n}}
%
\newcommand{\betareduce}{\leadsto_\beta}
\newcommand{\betareduceto}{\twoheadrightarrow_\beta}
\newcommand{\bit}{\mathtt{B}}
\newcommand{\bitone}{\mathtt{B}_1}
\newcommand{\bitzero}{\mathtt{B}_0}
%
\newcommand{\caseto}{\triangleright}
\newcommand{\cbind}[3]{\mathtt{bind} \ {#1} = {#2} \ \mathtt{in} \ {#3}}
\newcommand{\ccase}[4]{\mathtt{case} \; {#1} \; \mathtt{of} \; \cnil \caseto {#2}; \;
                                          \bitzero x_0 \caseto {#3}; \; \bitone x_1 \caseto {#4}}
\newcommand{\cequiv}{\simeq}
\newcommand{\cif}[3]{\mathtt{if} \; {#1} \; \mathtt{then} \; {#2} \; \mathtt{else} \; {#3}}
\newcommand{\ctrue}{\mathtt{true}}
\newcommand{\cfalse}{\mathtt{false}}
\newcommand{\cfix}{\mathtt{fix}\,}
\newcommand{\cfixin}[2]{\mathtt{fix}\; {#1}\; \mathtt{in} \; {#2}}
\newcommand{\cfold}{\mathtt{fold}\,}
\newcommand{\chd}{\mathtt{hd}}
\newcommand{\cinj}{\mathtt{inj}}
\newcommand{\clift}[1]{\lfloor{#1}\rfloor}
\newcommand{\clet}[3]{\mathtt{let} \ {#1} = {#2} \ \mathtt{in} \ {#3}}
\newcommand{\cletst}[3]{\mathtt{liftlet} \ {#1} = {#2} \ \mathtt{in} \ {#3}}
\newcommand{\cletrec}[3]{\mathtt{letrec} \ {#1} = {#2} \ \mathtt{in} \ {#3}}
\newcommand{\cnil}{\mathtt{nil}}
\newcommand{\conc}{\! \bullet \!}
\newcommand{\concat}{\pmb{\mathit{conc}}}
\newcommand{\context}{\mathbb{C}}
\newcommand{\cop}{\;\mathtt{op}\;}
\newcommand{\cpair}[1]{\langle{#1}\rangle}
\newcommand{\cproj}{\mathtt{proj}}
\newcommand{\crand}{\mathtt{rand}}
\newcommand{\crec}{\mathtt{rec}}
\newcommand{\croll}{\mathtt{roll}\,}
\newcommand{\crule}[1]{\mbox{{\textsl{#1}}}}
\newcommand{\cslrtype}{\Vdash}
\newcommand{\csucc}{\mathtt{S}}
\newcommand{\ctl}{\mathtt{tl}}
\newcommand{\cut}{\pmb{\mathit{cut}}}
\newcommand{\cval}{\mathtt{return}}
\newcommand{\cwhile}[2]{\mathtt{while}\; {#1} \; \mathtt{do} \; {#2} \; \mathtt{od}}
\newcommand{\cvalof}{\pmb{valueof}}
\newcommand{\cunfold}{\mathtt{unfold}\,}
\newcommand{\cunit}{\mathtt{unit}}
\newcommand{\cunroll}{\mathtt{unroll}\,}
\newcommand{\cupdate}{\pmb{update}}
\newcommand{\czero}{\mathtt{O}}
%
\newcommand{\deduces}{\vdash}
\newcommand{\defeq}{\stackrel{\mathrm{def}}{=}}
\newcommand{\delay}[1]{\left\lfloor{#1}\right\rfloor}
\newcommand{\distr}{\mathcal{D}}
%
\newcommand{\eagerreduce}{\xrightarrow{\mathrm{eager}}}
\newcommand{\Eagerreduce}{\stackrel{\mathrm{eager}}{\Reduce}}
\newcommand{\eqtest}{\stackrel{\mbox{\tiny ?}}{=}}
\newcommand{\Eval}[1]{{\left\llbracket{#1}\right\rrbracket}}
\newcommand{\evalcontext}{\mathbb{E}}
\newcommand{\Example}{\emgreen{\pmb{\textsc{Example}\;}}}
\newcommand{\expansion}[1]{\tilde{{#1}}}
%
\newcommand{\ford}{\mathrm{ord}}
%
\newcommand{\getval}{\leftarrow}
%
\newcommand{\head}{\pmb{\mathit{hd}}}
\newcommand{\termhole}{[\;]}
\newcommand{\hoaretriple}[3]{\{{#1}\} \; {#2} \; \{{#3}\}}
%
\newcommand{\imply}{\Rightarrow}
\newcommand{\indist}{\simeq}
\newcommand{\iteration}{\pmb{\mathit{it}}}
%
\newcommand{\lazyreduce}{\xrightarrow{\mathrm{lazy}}}
\newcommand{\Lazyreduce}{\lazyreduce\!\!\!\!\!\to}
\newcommand{\ldot}{\, . \, }
\newcommand{\linearto}{\multimap}
\newcommand{\logrel}{(\mathcal{R}_\tau)_{\mbox{$\tau$ type}}}
\newcommand{\lr}[1]{\mathcal{R}_{#1}}
\newcommand{\lub}{\sqcup}
%
\newcommand{\metalambda}{\underline{\lambda}}
\newcommand{\modal}{\square}
%
\newcommand{\nat}{\mathbb{N}}
\newcommand{\nReduce}{\nrightarrow\!\!\!\!\to}
\newcommand{\nreduceto}{\nrightsquigarrow}
\newcommand{\nullenv}{\varnothing}
%
\newcommand{\orc}{\mathcal{O}}
\newcommand{\oslrtype}{\Vdash}
%
\newcommand{\porder}{\sqsubseteq}
\newcommand{\prefix}{\pmb{\mathit{pref}}}
\newcommand{\proba}{\mathbf{Pr}}
%
\newcommand{\randlen}{\pmb{\mathit{rl}}}
\newcommand{\randnum}{\pmb{\mathit{rn}}}
\newcommand{\randstr}{\pmb{\mathit{rs}}}
\newcommand{\reduce}{\leadsto}
\newcommand{\reduceto}{\rightsquigarrow\hspace*{-6px}\shortrightarrow}
\newcommand{\Reduce}{\to\!\!\!\to}
\newcommand{\Reduceto}{\Downarrow}
%
\newcommand{\set}[1]{\{{#1}\}}
\newcommand{\setbits}{\mathbb{B}}
\newcommand{\slrtype}{\vdash}
\newcommand{\splt}{\pmb{\mathit{split}}}
\newcommand{\suffix}{\pmb{\mathit{suff}}}
\newcommand{\sysname}{\pmb{\mathcal{C}}}
%
\newcommand{\tail}{\pmb{\mathit{tl}}}
\newcommand{\tbits}{\mathsf{Bits}}
\newcommand{\tbool}{\mathsf{Bool}}
\newcommand{\tnat}{\mathsf{Nat}}
\newcommand{\tcomp}{\mathsf{T}}
\newcommand{\tlift}[1]{\lfloor{#1}\rfloor}
\newcommand{\tlist}{\mathsf{List}}
\newcommand{\tself}{\mathsf{self}}
\newcommand{\tstore}{\mathsf{Store}}
\newcommand{\tstream}{\mathsf{Stream}}
\newcommand{\tloc}{\mathsf{Loc}}
\newcommand{\ttcase}{\mathtt{case}}
\newcommand{\ttif}{\mathtt{if}}
\newcommand{\ttthen}{\mathtt{then}}
\newcommand{\ttelse}{\mathtt{else}}
\newcommand{\ttin}{\mathtt{in}}
\newcommand{\ttlet}{\mathtt{let}}
\newcommand{\ttof}{\mathtt{of}}
\newcommand{\tunit}{\mathsf{Unit}}
\newcommand{\Type}{\mathcal{T}\!\mathit{ype}}
%
\newcommand{\vbool}{\pmb{\mathbb{B}}}
\newcommand{\vff}{\mathfrak{f}\!\mathfrak{f}}
\newcommand{\vnat}{\pmb{\mathrm{N}}}
\newcommand{\vnil}{\epsilon}
\newcommand{\vtt}{\mathfrak{t}\!\mathfrak{t}}
%
\newcommand{\weakpre}{\pmb{wp}}




\pdsetup{lf={Theory and Practice of Functional Programming}}

\title{{\small \textsc{Theory and Practice of Functional Programming}}
\\[12pt]
{\Large \color{TitleColor} Type Inference}
}
\author{\small \textbf{Dr. ZHANG Yu}
  \\[2pt] \small Institute of Software, Chinese Academy of Sciences 
  \\[5pt] \includegraphics[height=.04\slideheight]{iscas.eps} 
  \qquad \qquad 
  \includegraphics[height=.04\slideheight]{cas.eps}
}
\date{\scriptsize Fall term, 2012 \qquad GUCAS, Beijing}

\newcommand{\RED}{\mathsf{RED}}

\begin{document}

\maketitle 

\begin{slide}[method=direct]{Principal types}
\begin{itemize}
\item A few type judgements:
\begin{smalltext}
\begin{eqnarray*} 
&& 
\vdash \lambda x : \tnat \ldot \lambda y : \tnat \to \tnat \ldot y x : 
\emblue{\tnat} \to (\emgreen{\tnat \to \tnat}) \to \empurple{\tnat}
\\ && 
\vdash \lambda x : \tbool \ldot \lambda y : \tbool \to \tnat \ldot y x : 
\emblue{\tbool} \to (\emgreen{\tbool \to \tnat}) \to \empurple{\tnat}
\\ && 
\vdash \lambda x : (\tnat \to \tbool) \ldot \lambda y : (\tnat \to \tbool) \to \tbool \to \tbool \ldot y x : 
\\ && \qquad
(\emblue{\tnat \to \tbool}) \to (\emgreen{(\tnat \to \tbool) \to \tbool \to \tbool}) \to \empurple{\tbool \to \tbool}
\end{eqnarray*}
\end{smalltext}
All types above are of the pattern $\emblue{X} \to (\emgreen{X \to Y}) \to \empurple{Y}$.
\item In Haskell, the system will give the following type to the term $\lambda x \ldot \lambda y \ldot y x$:
  %% # fun x -> fun y -> y x ;;
  %% - : 'a -> ('a -> 'b) -> 'b = <fun>
\vspace*{-1.8em}
\begin{smalltext}
\begin{verbatim}
    Prelude> :t \x -> \y -> y x
    \x -> \y -> y x :: t1 -> (t1 -> t) -> t
\end{verbatim}
\end{smalltext}
\item $X \to (X \to Y) \to Y$ is the most general type (or the \emblue{principal type}) for the term 
$\lambda x \ldot \lambda y \ldot y x$.
\\
Every typable lambda term in the simply typed lambda calculus has a principal type.
\end{itemize}
\end{slide}

\begin{slide}{Type templates}
\begin{itemize}
\item A \emblue{type template} is a simple type, built from \emblue{type variables} and possibly base types.
\[
T, S, \ldots ::= X \mid b \mid T \to S \mid T \times S \mid T + S
\]
where $X$ is type variable, $b$ ranges over a set of type constants.
\item A \emblue{type substitution} is a function from type variables to type templates: 
\vspace*{-0.5em}
  $$[X_1 \mapsto T_1, \ldots, X_n \mapsto T_n]  \vspace*{-0.5em}$$
We define the template $\sigma(T)$, obtained by applying the substitution $\sigma$ to templpate $T$, 
as follows:
\begin{smalltext}
\vspace*{-0.5em}
\begin{eqnarray*}
\sigma(X) & = & T, \mbox{ if $X \mapsto T \in \sigma$}
\\
\sigma (b) & = & b
\\
\sigma (T \to S) & = & \sigma(T) \to \sigma(S) 
\\
\sigma (T \times S) & = & \sigma(T) \times \sigma(S) 
\\[-2em]
\end{eqnarray*}
\end{smalltext}
A template $T$ is \emblue{more general} than $S$ if there exists a substitution $\sigma$ 
such that $\sigma(T) = S$, or $S$ instantiates $T$.
\end{itemize}
\end{slide}

\begin{slide}{Unification}
\begin{itemize}
\item Given an equation between type templates, a \emblue{unifier} is a type substitution
  that solves the equation. 
\begin{itemize}\small
\item $X \to (X \to Y) = (Y \to Z) \to W$ has a solution 
\[
\sigma = [ X \mapsto Y \to Z, W \mapsto (Y \to Z) \to Y ]
\vspace*{-1.5em}
\]
\item The equation $X \times (X \to Z) = (Z \to Y) \times Y$ does not have any solution.
\end{itemize}
\item A unifier of two sequences of type templates $\tilde{T} = T_1, \ldots , T_n$ and 
  $\tilde S = S_1, \ldots , S_n$, is a type substitution $\sigma$ s.t. 
  $\sigma(T_i) = \sigma(S_i)$ for all $i$.
\item A types substitution is a \emblue{most general unifier} of $T = S$ if it is a unifier 
and it is more general than any other unifier of $T =S$. \\
\begin{itemize} \small
\item A type substitution $\sigma$ is {\em more general} than $\sigma'$ if there exists another 
substitution $\rho$ such that $\rho \circ \sigma = \sigma'$.
\item Let $\sigma = [X_1 \mapsto T_1, \ldots , X_n \mapsto T_n]$. 
\emblue{$$\rho \circ \sigma = [X_1 \mapsto \rho(T_1), \ldots, X_n \mapsto \rho (T_n)] $$}
\end{itemize}
\end{itemize}
\end{slide}

\begin{slide}{The unification algorithm}
\begin{itemize}
\item \emblue{Unification} is the process of determining a most general unifier.
\vspace*{-0.5em}
\begin{itemize}\small
\item Input: two sequences of type templates $\tilde T$ and $\tilde S$.
\item Output: a most general unifier $\sigma$ if it exists.
\vspace*{-1em}
\end{itemize}
\item The most general unification algorithm \emred{$MGU$}:
\vspace*{-0.5em}
\begin{itemize}\small
\item $MGU(T; T) = \mathrm{id}$, if $T$ is a variable or a base type;
\item $MGU(X; S) = [ X \mapsto S ]$, if \empink{$X$ does not occur in $S$};
\item $MGU(T; Y) = [ Y \mapsto T ]$, if \empink{$Y$ does not occur in $T$};
\item $MGU(T_1 \to T_2; S_1 \to S_2) = MGU (T_1, T_2; S_1, S_2)$;
\item $MGU(T_1 \times T_2; S_1 \times S_2) = MGU (T_1, T_2; S_1, S_2)$;
\item $MGU(T_1 + T_2; S_1 + S_2) = MGU (T_1, T_2; S_1, S_2)$;
\item $MGU(T; S)$ fails in all other cases;
\item $MGU(T, \tilde T; S, \tilde S) = \sigma' \circ \sigma$, where $\sigma = MGU(\tilde T; \tilde S)$ 
  and $\sigma' = MGU(\sigma(T); \sigma(S))$.
\vspace*{-0.5em}
\end{itemize}
\item If $MGU(\tilde T; \tilde S) = \sigma$, then $\sigma$ is a most general unifier of $\tilde T$ 
  and $\tilde S$.
\end{itemize}
\end{slide}

\begin{slide}{The type inference algorithm}
The type inference algorithm $TI$: \emblue{$TI(\Gamma \vdash e: {T})$} is a most general substitution $\sigma$ 
such that \emred{$\sigma(\Gamma) \vdash e : \sigma(T)$ holds} ($TI$ fails if no such $\sigma$).
\vspace*{-0.5em}
\begin{itemize}\small
\item[--] $TI(x_1 : T_1, \ldots , x_n : T_n \vdash x_i : T) = MGU(T_i; T)$. \vspace*{-0.5em}
\item[--] $TI(\Gamma \vdash e_1 + e_2 : T) = \sigma_2 \circ \sigma_1 \circ \sigma$, where \\
  \begin{minipage}{.5\textwidth}
  \begin{itemize}
  \item $\sigma = MGU(\tnat; T)$
  \item $\sigma_1 = TI(\sigma(\Gamma) \vdash e_1 : \tnat)$
  \item $\sigma_2 = TI(\sigma_1 \circ \sigma(\Gamma) \vdash e_1 : \tnat)$
  \end{itemize}
  \end{minipage}
  \hfill
  \begin{minipage}{.3\textwidth}
    $$\prooftree \Gamma \vdash e_1 : \tnat \quad \Gamma \vdash e_2: \tnat \justifies 
    \Gamma \vdash e_1 + e_2 : \tnat \endprooftree$$
  \end{minipage} 
  \vspace*{-1em}
\item[--] $TI(\Gamma \vdash e_1 e_2 : T) = \sigma_2 \circ \sigma_1$, where \\
  \begin{minipage}{.5\textwidth}
  \begin{itemize}
  \item $\sigma_1 = TI(\Gamma \vdash e_1 : X \to T)$
  \item $\sigma_2 = TI(\emblue{\sigma_1(\Gamma)} \vdash e_2 : \emblue{\sigma_1(X)})$ 
  \item $X$ is a fresh type variable
  \end{itemize}
  \end{minipage}
  \hfill
  \begin{minipage}{.3\textwidth}
    $$\prooftree \Gamma \vdash e_1 : \tau \to \tau' \quad \Gamma \vdash e_2: \tau \justifies 
    \Gamma \vdash e_1 \ e_2 : \tau' \endprooftree$$ 
  \end{minipage}
  \vspace*{-1em}
\item[--] $TI(\Gamma \vdash \lambda x: S \ldot e : T) = \sigma_2 \circ \sigma_1$,  where \\
  \begin{minipage}{.6\textwidth}
  \begin{itemize}
  \item $\sigma_1 = MGU(T; S \to X)$
  \item $\sigma_2 = TI(\sigma_1(\Gamma), x: \sigma_1(S) \vdash e : \sigma_1(X))$
  \item $X$ is a fresh variable
  \end{itemize}
  \end{minipage}
  \hfill
  \begin{minipage}{.3\textwidth}
    $$\prooftree \Gamma, x : \tau \vdash e : \tau' 
    \justifies \Gamma \vdash \lambda x \ldot e : \tau \to \tau' \endprooftree$$
  \end{minipage}
\end{itemize}
\end{slide}

\begin{slide}{The type inference algorithm}
\begin{itemize}\small
\item[--] $TI(\Gamma \vdash \cpair{e_1, e_2} : T) = \sigma_2 \circ \sigma_1 \circ \sigma$, where 
  \begin{minipage}{.5\textwidth}
  \begin{itemize}
  \item $\sigma = MGU(T; X \times Y)$
  \item $\sigma_1 = TI(\sigma(\Gamma) \vdash e_1 : \sigma(X))$
  \item $\sigma_2 = TI(\sigma_1 \circ \sigma (\Gamma) \vdash e_2 : \sigma_1 \circ \sigma(Y))$ 
  \item  $X, Y$ are fresh type variables
  \end{itemize}
  \end{minipage}
  \begin{minipage}{.3\textwidth}
    $$\prooftree \Gamma \vdash e_1 : \tau_1 \quad \Gamma \vdash e_2 : \tau_2
    \justifies \Gamma \vdash \cpair{e_1, e_2} : \tau_1 \times \tau_2 \endprooftree$$
  \end{minipage}
  \vspace*{-1em}
\item[--] $TI(\Gamma \vdash \cproj_1 e : T) = TI(\Gamma \vdash e : T \times X)$, where 
  \begin{minipage}{.5\textwidth}
  \begin{itemize}
  \item $X$ is a fresh type variable
  \item $TI(\Gamma \vdash \cproj_2 e : T)$ is similar.
  \end{itemize}
  \end{minipage}
  \begin{minipage}{.3\textwidth}
    $$\prooftree \Gamma \vdash e : \tau_1 \times \tau_2
    \justifies \Gamma \vdash \cproj_i \ e : \tau_i \endprooftree$$
  \end{minipage}
  \vspace*{-1em}
\item[--] $TI(\Gamma \vdash \cinj_1 e : T) = TI(\Gamma \vdash e : T \times X)$, where 
  \begin{minipage}{.5\textwidth}
  \begin{itemize}
  \item $\sigma = MGU(X + Y; T)$ 
  \item $\sigma_1 = TI(\sigma{\Gamma} \vdash e : X$
  \item $X, Y$ is a fresh type variable
  \item $TI(\Gamma \vdash \cinj_2 e : T)$ is similar.
  \end{itemize}
  \end{minipage}
  \begin{minipage}{.3\textwidth}
    $$\prooftree \Gamma \vdash e : \tau_i
    \justifies \Gamma \vdash \cinj_i \ e : \tau_1 + \tau_2 \endprooftree$$
  \end{minipage}
\end{itemize}
\end{slide}

\begin{slide}{The type inference algorithm}
\begin{itemize}\small
\item[--] $TI(\Gamma \vdash \mathtt{case} \; e \; \mathtt{of} \; \cinj_1 x_1 \; \mathtt{in} \; e_1 ; \cinj_2 x_2 \; \mathtt{in} \; e_2 : T) = \sigma_2 \circ \sigma_1 \circ \sigma$, where 
  \begin{itemize}
  \item $\sigma = TI(\Gamma \vdash e :X + Y)$
  \item $\sigma_1 = TI(\sigma(\Gamma), x_1 : \sigma(X) \vdash e_1 : \sigma(T))$
  \item $\sigma_2 = TI(\sigma_1 \circ \sigma (\Gamma), x_2 : \sigma_1 \circ \sigma (Y) \vdash e_2 : \sigma_1 \circ \sigma(Y))$ 
  \item  $X, Y$ are fresh type variables
  \end{itemize}
$$
\prooftree
\Gamma, x_1 : \tau_1 \vdash e_1 : \tau 
\quad 
\Gamma, x_2 : \tau_2 \vdash e_2 : \tau 
\quad 
\Gamma \vdash e : \tau_1 + \tau_2
\justifies
\Gamma \vdash 
\mathtt{case} \; e \; \mathtt{of} \; \cinj_1 x_1 \; \mathtt{in} \; e_1 ; \cinj_2 x_2 \; \mathtt{in} \; e_2 : \tau
\endprooftree
$$
\item[--] $TI(\Gamma \vdash \cfix e : T) = TI(\Gamma \vdash e : T \to T)$.
$$
  \prooftree \Gamma \vdash e : \tau \to \tau 
  \justifies \Gamma \vdash \cfix \ e : \tau \endprooftree
$$
\end{itemize}
%% \item $TI(\Gamma \vdash e : T)$ returns a most general substitution $\sigma$ such that
%%   $\sigma(\Gamma) \vdash e : \sigma(T)$ holds. It fails if no such substitution exists.
In real applications, type inference only takes two arguments $\Gamma$ and $e$, 
and we do inference by $TI(\Gamma \vdash e : \emred{X})$, with $X$ a fresh type variable.
\end{slide}


\end{document}
 

